\section{Related Work}
\label{sec:related}

\para{Prompt sensitivity and robustness.}
\citet{sclar2023quantifying} introduced \textsc{FormatSpread}, demonstrating that meaning-preserving prompt format changes cause up to 76 accuracy point swings across 53 tasks. Their formal grammar over prompt formats (separators, spacing, casing) reveals a highly non-monotonic performance landscape.
\citet{mizrahi2024state} conducted a large-scale study (6.5M instances, 20 models) confirming that prompt paraphrasing effects are severe: 10 of 25 tasks show weak ranking agreement across prompts, and edit distance does not predict degradation.
\citet{zhu2023promptbench} proposed a benchmark covering character-, word-, sentence-, and semantic-level perturbations.
These studies thoroughly document prompt fragility but do not provide mechanistic explanations or predictive tools. Our work complements this line of research by connecting internal model geometry to prompt-level failure prediction.

\para{Attention geometry and model robustness.}
\citet{park2024geometry} treat attention maps as weighted graphs and compute Ollivier-Ricci curvature (ORC), establishing that higher positive curvature correlates with more robust transformers.
They show that attention variance serves as a computationally efficient proxy for ORC, and that ORC distributions shift toward positive values through the forward pass.
However, their analysis operates at the model level during training, not at the prompt level during inference.
\citet{piotrowski2025constrained} show that transformers implement constrained Bayesian belief updating, with internal representations forming fractal structures in the belief simplex.
Their finding that architectural constraints warp optimal inference geometry provides theoretical grounding for why attention geometry should vary with prompt format.
We extend these model-level geometric insights to prompt-level failure prediction at inference time.

\para{Attention head analysis and circuit discovery.}
\citet{wang2022interpretability} identified the complete circuit for indirect object identification (IOI) in \gptsmall, including name movers, induction heads, and S-inhibition heads.
\citet{olsson2022context} identified induction heads as a key mechanism for in-context learning.
\citet{conmy2023towards} automated circuit discovery with ACDC, systematically identifying task-relevant edges in the computational graph.
\citet{bussman2025beyond} introduced SVD-based circuit analysis, decomposing attention heads into singular vector components.
These works establish that attention heads form interpretable, task-specific circuits, motivating our extraction of geometric features from these circuits.

\para{Circuit geometry and stability.}
\citet{suarez2026circuit} propose that circuit discovery and activation steering are dual operations on the same geometric structure, finding that negative-valence steering is fragile---directly demonstrating geometric regions of instability.
\citet{bali2026quantifying} measure attention head stability across random seeds, finding a U-shaped stability curve where mid-layer heads are least stable yet most functionally important.
They also show that stability degrades with prompt length and that query-weight Frobenius norm correlates negatively with stability.
Our work is closest to these studies in spirit, but we focus on a different question: whether geometric features predict \emph{prompt-level} performance rather than \emph{model-level} stability.

\para{Representation geometry.}
\citet{park2024linear} analyze when linear representations emerge in LLMs, while \citet{engels2024not} show that some features are represented as multi-dimensional, non-linear structures.
These findings suggest that attention circuit geometry may have complex manifold structure relevant to prompt sensitivity, supporting our use of multiple geometric feature types (entropy, singular values, variance) rather than a single geometric measure.
